\chapter{Derivative Validation}
\label{ch:bf-derivative-validation}

Derivative validation is a ladder.  Passing one rung is useful evidence, but it
does not certify the whole industrial target.

\section{Validation Ladder}
\label{sec:bf-derivative-validation-ladder}

BayesFilter derivative providers should pass:
\begin{enumerate}[label=(\roman*)]
  \item shape and finite-value checks;
  \item finite-difference checks on scalar or tiny LGSSMs;
  \item autodiff parity on small models where tape Hessians are feasible;
  \item backend parity among covariance, solve, Cholesky, and QR variants;
  \item factor reconstruction and solve-residual diagnostics;
  \item compiled/eager parity for TensorFlow backends;
  \item stress tests for low noise, collinearity, near rank loss, and masks.
\end{enumerate}

MacroFinance currently supplies examples of these tests, including QR
square-root parity against autodiff and finite differences, solve backend
agreement with the existing differentiated backend, TensorFlow eager/compiled
parity, and large-scale backend ladder diagnostics.  BayesFilter should import
the testing discipline before importing any claim of production readiness.

\section{What Is Being Compared}
\label{sec:bf-validation-comparison-target}

Every derivative test must name the scalar or tensor being compared.  For a
sampler-coordinate posterior target,
\[
  \mathcal{L}(u)
  =
  \ell(T(u);y_{1:T})
  + \log p(T(u))
  + \log\left|\det DT(u)\right|,
\]
the score test compares
\[
  \hat g_i(u)
  \quad\hbox{to}\quad
  \frac{\mathcal{L}(u+h e_i)-\mathcal{L}(u-h e_i)}{2h}
\]
or to an autodiff score of that same scalar.  The Hessian test compares
\[
  \hat H_{ij}(u)
  \quad\hbox{to}\quad
  \frac{\hat g_i(u+h e_j)-\hat g_i(u-h e_j)}{2h}
\]
or to a tape Hessian of the same scalar.  If the reported analytic derivative is
likelihood-only but the finite difference includes the prior or transform
Jacobian, the test is invalid.  Conversely, if a chapter validates
\(\nabla_\theta\ell\) in model-parameter coordinates, that result should be read
as a component-level backend check rather than as the complete HMC target
gradient unless the prior and Jacobian terms are added consistently.  If the
value path contains jitter, floors, masks, transformed parameters, or invalid-
region fallback values, the comparison value must contain the same choices.

\section{Recommended Rungs}
\label{sec:bf-validation-recommended-rungs}

The derivative validation ladder should be run in this order:
\begin{enumerate}[label=(\roman*)]
  \item \textbf{Finite values and shapes.}
    Check that value, score, Hessian, state moments, factors, provider tensors,
    and diagnostics are finite and have declared dimensions.
  \item \textbf{Symmetry and admissibility.}
    Check that covariance-like objects are symmetric, positive
    semidefinite/definite under their declared policy, and that Hessians are
    symmetric within tolerance.
  \item \textbf{Block finite differences.}
    Validate structural derivatives of $c,F,Q,a,H,R,m_0,P_0$ before running a
    long filter scan.
  \item \textbf{Likelihood finite differences.}
    Compare score and Hessian of the complete scalar target on tiny systems.
  \item \textbf{Autodiff parity.}
    Use taped gradients and Hessians as small-system oracles where nested tapes
    are feasible.
  \item \textbf{Backend parity.}
    Compare covariance, solve, Cholesky, QR, and SVD/spectral backends on
    systems where they represent the same law.
  \item \textbf{Factor reconstruction.}
    Check
    \eqref{eq:bf-factor-reconstruction-first}--\eqref{eq:bf-factor-reconstruction-second}
    for prediction, innovation, and filtered covariance factors.
  \item \textbf{Mask tests.}
    Test dense all-observed masks, sparse deterministic masks, and all-missing
    periods using the convention in Section~\ref{sec:bf-score-masks}.
  \item \textbf{Spectral and branch stress tests.}
    Vary spectral gaps, Cholesky pivots, QR diagonals, rank thresholds, floors,
    jitter, and fallback branches while checking diagnostics.
  \item \textbf{Eager/compiled parity.}
    Compare values and derivatives under eager execution, compiled execution,
    and XLA/JIT paths when those paths are intended for production.
  \item \textbf{Sampler smoke tests.}
    Run short HMC/NUTS checks only after deterministic value-and-gradient tests
    pass, and interpret them as sampler usability evidence, not proof of the
    derivative formulas.
\end{enumerate}

\section{Finite-Difference Tolerances}
\label{sec:bf-validation-finite-difference-tolerances}

Finite differences are not exact tests; they are numerical experiments.  The
step size should be selected relative to the transformed parameter scale, for
example
\[
  h_i = \eta(1+|\psi_i|),
\]
with several $\eta$ values checked when possible.  A central difference should
show a stable error plateau before the result is accepted.  Parameters near
constraints, branch boundaries, or spectral floors should be tested separately
from smooth interior points, because the expected derivative may be a branch
derivative rather than a smooth derivative of the unregularized model.

\section{Backend-Specific Gates}
\label{sec:bf-validation-backend-gates}

Each backend has a distinct failure mode:
\begin{itemize}
  \item covariance form can lose symmetry or positive semidefiniteness under
    finite arithmetic;
  \item solve form can hide ill conditioning unless solve residuals and
    log-determinant policies are reported;
  \item Cholesky form depends on positive pivots and any jitter policy;
  \item QR form depends on full column rank, diagonal sign convention, and
    pivot policy;
  \item SVD/spectral form depends on gap, rank, floor, sign/order, and fallback
    policy;
  \item sigma-point form additionally depends on fixed point rules, weights,
    nonlinear map derivatives, and moment reconstruction.
\end{itemize}
Backend parity is meaningful only when the compared backends target the same
law.  A covariance backend using $P$ and an SVD backend using a floored $P_+$ are
not required to have identical derivatives unless $P_+=P$ on that test case.

\section{Validation Artifact}
\label{sec:bf-validation-artifact}

Every promoted derivative backend should leave an artifact containing:
\begin{itemize}
  \item model name, parameter vector, transform policy, and reference point;
  \item backend name and value branch;
  \item tests run, tolerances, and maximum absolute/relative errors;
  \item finite-value, finite-gradient, and finite-Hessian status;
  \item branch diagnostics: jitter, PSD projection, spectral floors, pivots,
    rank truncation, sign/order corrections, and fallbacks;
  \item interpretation: which claim is supported and which regimes remain
    untested.
\end{itemize}
This artifact is the bridge between a mathematical derivation and an HMC claim.
Without it, phrases such as ``analytic Hessian'', ``SVD sigma-point gradient'',
or ``HMC-ready backend'' are too ambiguous to audit.

\section{Audit Record}
\label{sec:bf-derivative-audit-record}

Each derivative chapter should record the source labels, code files, and tests
that support it.  If a tool audit is inconclusive, the monograph should say so.
An inconclusive symbolic audit can still be valuable: it identifies which
guards, shapes, or noncommutative matrix steps need manual proof before a
formula is promoted to peer-review-ready status.

For the analytic-derivative chapter set, the current audit record is:
\begin{itemize}
  \item MacroFinance analytic derivative labels provide the source formulas for
    score, Hessian, structural provider, Cholesky, and QR contracts;
  \item MathDevMCP is useful for label lookup and local-context provenance, but
    it does not replace manual proof of noncommutative matrix calculus steps;
  \item ResearchAssistant is local/offline and has no current summary hits for
    the exact SVD/eigen derivative-validation cluster used here;
  \item production readiness remains a test claim, not a prose claim.
\end{itemize}
